%% title: Asymptotics of NLS and Hypothesis Testing with the GNR
\documentclass[11pt]{article}
\usepackage{metricsstyle}

\begin{document}

\lessonheader{12}{Nonlinear least squares: consistency, asymptotic normality and hypothesis testing with the GNR}{October 28, 2026}

\runin{Previously} The Gauss--Newton regression (GNR) and Taylor series (Lesson 11).
\runin{Announcements} November 9: exam assignment. November 23: Assignment 2.

\section{Recap: linear and nonlinear regression}

\topic{The linear model}
\[ y = X\beta + u, \qquad u \sim (0,\sigma^{2}I_{n}), \qquad
   \E(x_{t}u_{t}) = 0, \]
\[ \operatorname*{plim}\frac{X'X}{n} = \Sigma_{XX}, \qquad
   \rank(\Sigma_{XX}) = K. \]
Under the law of large numbers and central limit theorem conditions used in Lesson 10,
\[ b_{n} = (X'X)^{-1}X'y, \qquad
   \sqrt{n}\,(b_{n}-\beta) \ \xrightarrow{\ d\ }\ N(0,\sigma^{2}\Sigma_{XX}^{-1}). \]

\topic{The nonlinear model}
\[ y = x(\beta_{0}) + u, \qquad u \sim (0,\sigma^{2}I_{n}), \]
where $\beta_{0}$ is the true parameter vector and $x(\beta)$ is the $n\times 1$
vector of regression functions. Define the $n\times K$ derivative matrix
\[ X(\beta) = \Bigl[\frac{\partial x_{t}(\beta)}{\partial\beta_{i}}\Bigr]_{t,i},
   \qquad t=1,\ldots,n,\quad i=1,\ldots,K. \]
As in Lesson 11, lower-case $x(\beta)$ is the regression function; capital
$X(\beta)$ is its derivative matrix. If $X_{t}(\beta)$ is its $t$-th row, the
orthogonality condition is
\[ \E\bigl(X_{t}(\beta_{0})'u_{t}\bigr)=0
   \quad\Longrightarrow\quad \E\bigl(X(\beta_{0})'u\bigr)=0. \]
The NLS estimator minimizes the sum of squared residuals:
\[ \hat\beta = \operatorname*{arg\,min}_{\beta}\operatorname{SSR}(\beta),
   \qquad \operatorname{SSR}(\beta)=(y-x(\beta))'(y-x(\beta)). \]

\Needspace{14\baselineskip}
\topic{Notation and the rank condition}
\[ \hat x=x(\hat\beta), \qquad x_{0}=x(\beta_{0}), \qquad
   \hat X=X(\hat\beta), \qquad X_{0}=X(\beta_{0}), \]
\[ \operatorname*{plim}\frac{X_{0}'X_{0}}{n}=\Sigma_{XX},
   \qquad \Sigma_{XX}\text{ finite and positive definite}. \]
This is the local rank condition. Consistency also requires identification of
$\beta_{0}$ as the unique minimizer of the population objective and suitable
uniform convergence. The derivations below assume these regularity conditions,
smoothness, and the required LLN and CLT. Conditional mean zero,
$\E(u_{t}\mid\text{regressors})=0$, together with conditional homoskedasticity
and suitable sampling conditions, is a sufficient setting for the moment
conditions used below.

\section{Consistency of the residual variance estimate}

\topic{Expand the sum of squared residuals}
Since $y=x_{0}+u$,
\begin{align*}
 \operatorname{SSR}(\hat\beta)
 &= (y-\hat x)'(y-\hat x)\\
 &= y'y-2y'\hat x+\hat x'\hat x\\
 &= (x_{0}+u)'(x_{0}+u)-2(x_{0}+u)'\hat x+\hat x'\hat x\\
 &= x_{0}'x_{0}+2u'x_{0}+u'u-2x_{0}'\hat x-2u'\hat x+\hat x'\hat x\\
 &= u'u+2u'(x_{0}-\hat x)+(x_{0}-\hat x)'(x_{0}-\hat x).
\end{align*}

\topic{Taylor approximation around $\beta_{0}$}
Write $\delta=\hat\beta-\beta_{0}$. A first-order expansion gives
\[ \hat x-x_{0}=X_{0}\delta+r_{n}, \]
where $r_{n}$ collects the higher-order terms. Under consistency and the
regularity conditions, their contribution to $\operatorname{SSR}/n$ is $o_{p}(1)$.
Substituting into the expansion,
\[ \frac{\operatorname{SSR}(\hat\beta)}{n}
   =\frac{u'u}{n}-2\frac{u'X_{0}}{n}\delta
     +\delta'\frac{X_{0}'X_{0}}{n}\delta+o_{p}(1). \]
Now use
\[ \frac{u'u}{n}\xrightarrow{\ p\ }\sigma^{2}, \qquad
   \frac{X_{0}'u}{n}\xrightarrow{\ p\ }0, \qquad
   \frac{X_{0}'X_{0}}{n}\xrightarrow{\ p\ }\Sigma_{XX}, \qquad
   \delta\xrightarrow{\ p\ }0. \]
Both the cross term and the quadratic term vanish. At the true parameter,
\[ \operatorname{SSR}(\beta_{0})=(y-x_{0})'(y-x_{0})=u'u. \]
Hence
\[ \boxed{\frac{\operatorname{SSR}(\hat\beta)}{n}\xrightarrow{\ p\ }\sigma^{2}},
   \qquad \frac{\operatorname{SSR}(\beta_{0})}{n}\xrightarrow{\ p\ }\sigma^{2}. \]
This establishes consistency of the residual variance estimate, given
consistency of $\hat\beta$; it does not itself prove consistency of $\hat\beta$.
For fixed $K$, the degrees-of-freedom adjustment has the same limit:
\[ s^{2}=\frac{\operatorname{SSR}(\hat\beta)}{n-K}\xrightarrow{\ p\ }\sigma^{2}. \]

\Needspace{12\baselineskip}
\section{Asymptotic normality of NLS}

\topic{First-order conditions and the derivative}
The NLS first-order conditions are
\[ g_{n}(\hat\beta)=\hat X'(y-\hat x)=0, \qquad
   g_{n}(\beta)=X(\beta)'(y-x(\beta)). \]
Let $H_{t}(\beta)=\partial^{2}x_{t}(\beta)/\partial\beta\,\partial\beta'$,
a $K\times K$ matrix. At $\beta_{0}$,
\[ g_{n}(\beta_{0})=X_{0}'u, \qquad
   \frac{\partial g_{n}(\beta_{0})}{\partial\beta'}
      =\sum_{t=1}^{n}u_{t}H_{t}(\beta_{0})-X_{0}'X_{0}. \]
Assume the second-derivative moment and LLN conditions
\[ \E\bigl(u_{t}H_{t}(\beta_{0})\bigr)=0, \qquad
   \frac{1}{n}\sum_{t=1}^{n}u_{t}H_{t}(\beta_{0})\xrightarrow{\ p\ }0. \]


\topic{Taylor expansion and scaling by $\sqrt{n}$}
Under the usual local remainder conditions, expand the normalized FOCs:
\[ 0=\frac{X_{0}'u}{\sqrt{n}}
   +\left[\frac{1}{n}\sum_{t=1}^{n}u_{t}H_{t}(\beta_{0})
          -\frac{X_{0}'X_{0}}{n}\right]\sqrt{n}(\hat\beta-\beta_{0})+o_{p}(1). \]
The bracket converges to $-\Sigma_{XX}$, giving
\[ \sqrt{n}(\hat\beta-\beta_{0})
   =\Sigma_{XX}^{-1}\frac{X_{0}'u}{\sqrt{n}}+o_{p}(1). \]

\topic{Distribution of $X_{0}'u/\sqrt{n}$}
For fixed regressors, under $\Var(u)=\sigma^{2}I_{n}$,
\[ \E\Bigl(\frac{X_{0}'u}{\sqrt{n}}\Bigr)=0, \qquad
   \Var\Bigl(\frac{X_{0}'u}{\sqrt{n}}\Bigr)
      =\frac{X_{0}'(\sigma^{2}I_{n})X_{0}}{n}
      =\sigma^{2}\frac{X_{0}'X_{0}}{n}\longrightarrow\sigma^{2}\Sigma_{XX}. \]
For stochastic regressors use the corresponding conditional assumptions. By the CLT,
\[ \frac{X_{0}'u}{\sqrt{n}}\xrightarrow{\ d\ }N(0,\sigma^{2}\Sigma_{XX}). \]

\topic{Putting the two pieces together}
\begin{align*}
 \sqrt{n}(\hat\beta-\beta_{0})
 &\xrightarrow{\ d\ }\Sigma_{XX}^{-1}N(0,\sigma^{2}\Sigma_{XX})\\
 &=N\bigl(0,\sigma^{2}\Sigma_{XX}^{-1}\Sigma_{XX}\Sigma_{XX}^{-1}\bigr)
 =N(0,\sigma^{2}\Sigma_{XX}^{-1}).
\end{align*}
Thus
\[ \boxed{\sqrt{n}(\hat\beta-\beta_{0})\xrightarrow{\ d\ }
   N(0,\sigma^{2}\Sigma_{XX}^{-1})}, \qquad
   \Sigma_{XX}=\operatorname*{plim}\frac{X_{0}'X_{0}}{n}. \]
This has the same form as the OLS result: the derivative matrix $X(\beta_{0})$
replaces the regressor matrix $X$.

\section{The GNR and the estimated covariance matrix}

At the unrestricted NLS estimate, the GNR is
\[ y-\hat x=\hat Xb+\text{errors}. \]
Its OLS coefficient vector is zero:
\[ \hat b=(\hat X'\hat X)^{-1}
   \underbrace{\hat X'(y-\hat x)}_{=\,0\text{ by the FOC}}=0. \]
Although its coefficients are zero, its usual OLS covariance formula provides
the estimated covariance matrix for the NLS estimator:
\[ \widehat{\Var}(\hat\beta)=s^{2}(\hat X'\hat X)^{-1},
   \qquad s^{2}=\frac{(y-\hat x)'(y-\hat x)}{n-K}. \]
Consistency and continuity imply $\hat X'\hat X/n\xrightarrow{p}\Sigma_{XX}$, so
\[ n\,\widehat{\Var}(\hat\beta)
   =s^{2}\Bigl(\frac{\hat X'\hat X}{n}\Bigr)^{-1}
      \xrightarrow{\ p\ }\sigma^{2}\Sigma_{XX}^{-1}. \]
The factor $n$ distinguishes the covariance estimate for $\hat\beta$ from the
limiting covariance matrix of $\sqrt{n}(\hat\beta-\beta_{0})$.

\section{Hypothesis testing with the restricted GNR}

\topic{Partition the parameters}
\[ y=x(\beta_{1},\beta_{2})+u, \qquad
   \beta_{1}:K_{1}\times1,\quad \beta_{2}:K_{2}\times1,\quad K=K_{1}+K_{2}. \]
We want to test the $K_{2}$ restrictions
\[ H_{0}:\beta_{2}=0. \]
Here $\beta_{1}$ and $\beta_{2}$ denote parameter \emph{blocks}; the subscripts
in the example below denote individual coefficients.

\Needspace{16\baselineskip}
\topic{Example: restrictions that make the model linear}
\[ y_{t}=\beta_{1}x_{1t}+\beta_{2}x_{2t}^{\beta_{3}}+u_{t},
   \qquad x_{2t}>0. \]
\begin{flatlist}
 \item $H_{0}:\beta_{3}=0$ gives $y_{t}=\beta_{1}x_{1t}+\beta_{2}+u_{t}$:
       the power term becomes a constant.
 \item $H_{0}:\beta_{3}=1$ gives $y_{t}=\beta_{1}x_{1t}+\beta_{2}x_{2t}+u_{t}$.
 \item $H_{0}:\beta_{2}=0$ removes the power term, but then $\beta_{3}$ is
       unidentified. The regular full-rank testing argument below does not
       apply directly to that null.
\end{flatlist}
A restriction such as $\beta_{3}=1$ can be written as a zero restriction after
reparameterizing with $\gamma=\beta_{3}-1$.

\topic{Partition the derivative matrix}
\[ X(\beta)=\begin{bmatrix}X_{1}(\beta)&X_{2}(\beta)\end{bmatrix}, \qquad
   X_{j}(\beta)=\frac{\partial x(\beta)}{\partial\beta_{j}'},\quad j=1,2. \]
The artificial regression is
\[ y-x(\beta)=X_{1}(\beta)b_{1}+X_{2}(\beta)b_{2}+\text{errors}. \]
At $\hat\beta$, both coefficient blocks are zero by the unrestricted FOCs.
For a score (LM) test, evaluate the GNR at the \emph{restricted} estimate instead.

\topic{Evaluate at the restricted estimate}
Let $\tilde\beta=(\tilde\beta_{1}',0')'$ minimize SSR under $H_{0}$, and write
\[ \tilde u=y-x(\tilde\beta), \qquad
   \tilde X_{1}=X_{1}(\tilde\beta), \qquad \tilde X_{2}=X_{2}(\tilde\beta). \]
The restricted FOCs give $\tilde X_{1}'\tilde u=0$, while
$\tilde X_{2}'\tilde u$ need not be zero. Run
\[ \tilde u=\tilde X_{1}b_{1}+\tilde X_{2}b_{2}+\text{errors}. \]
The coefficient on $\tilde X_{1}$ is zero if that block is used alone. In the
\emph{joint} regression, however, its coefficient need not be zero, because
$\tilde X_{1}$ and $\tilde X_{2}$ need not be orthogonal. With
\[ M_{1}=I_{n}-\tilde X_{1}(\tilde X_{1}'\tilde X_{1})^{-1}\tilde X_{1}', \]
FWL gives
\[ \hat b_{2}=(\tilde X_{2}'M_{1}\tilde X_{2})^{-1}\tilde X_{2}'M_{1}\tilde u,
   \qquad M_{1}\tilde u=\tilde u. \]
Test the block $b_{2}=0$ in this restricted GNR. Under the regularity and
homoskedasticity assumptions, the usual auxiliary-regression $F$ statistic
satisfies $K_{2}F\xrightarrow{d}\chi^{2}_{K_{2}}$ under $H_{0}$; this is an
asymptotic test, rather than an exact finite-sample $F$ result.

\topic{Intuition}
If $H_{0}$ is true, $\tilde\beta$ and $\hat\beta$ approach the same true value,
and the restricted GNR coefficients should be small. A large, significant
$\hat b_{2}$ indicates that moving away from the restriction can improve the fit.
The unrestricted GNR still supplies covariance estimates for Wald tests; its
zero coefficients only prevent this particular score test from being based there.

\clearpage
\begin{transcriptnotes}
 \item Source: the five-page typed export \emph{Lesson 12.pdf}, dated October 28,
       2026. The title block, numbered sections, bold topic headings, prime
       transposes, $K$ notation and displayed derivations follow Lessons 10--11
       and the shared \texttt{metricsstyle} package. The contents list and coloured
       boxes are omitted; announcements are retained as recorded in the source.
 \item Pages 1--2: $X_{t}(\beta_{0})$ is defined as a row, so its transpose is
       included in the per-observation orthogonality condition. The model is
       evaluated at $\beta_{0}$. The NLS objective and the distinction between
       local rank and global identification are made explicit. Regularity,
       conditional moment and sampling conditions are stated to qualify the
       source's asymptotic calculations.
 \item Pages 2--3: the section headed ``Consistency'' assumes
       $\hat\beta\xrightarrow{p}\beta_{0}$; it establishes consistency of
       $\operatorname{SSR}(\hat\beta)/n$, not consistency of $\hat\beta$ itself.
       The heading is clarified. The Taylor remainder is retained as $r_{n}$,
       with its normalized contribution written $o_{p}(1)$; the source's
       approximate substitution is not treated as an exact equality.
       The definition and consistency of $s^{2}$ are added.
 \item Pages 3--4: the ambiguous derivative $\partial X'/\partial\beta\,u$
       is written $\sum_{t}u_{t}H_{t}(\beta_{0})$. Its mean-zero condition and an
       LLN are required explicitly; $\E(X_{0}'u)=0$ alone does not imply either.
       Taylor remainders are stated for the normalized FOCs as $o_{p}(1)$.
       The variance calculation is qualified for fixed or conditionally
       exogenous regressors, and the CLT is kept as a separate requirement.
 \item Page 4: the ordinary GNR covariance formula is distinguished from its
       $n$-scaled limit, the asymptotic covariance of $\sqrt{n}(\hat\beta-\beta_{0})$.
       It estimates the NLS covariance asymptotically, rather than being an
       exact finite-sample variance identity.
 \item Page 5: \textbf{Corrected:} $\beta_{3}=0$ makes
       $\beta_{2}x_{2t}^{\beta_{3}}=\beta_{2}$; it does not remove the term and
       is not equivalent to $\beta_{2}=0$. The latter null leaves $\beta_{3}$
       unidentified. \textbf{Corrected:} the restricted FOC makes
       $\tilde X_{1}'\tilde u=0$, but does not force the coefficient on
       $\tilde X_{1}$ to zero in the joint GNR. The FWL formula explains this.
       The block derivative notation, reparameterization, asymptotic $F$/LM
       calibration and distinction from unrestricted Wald inference are added
       to clarify how the proposed test works.
 \item Original page map: page 1 title, announcements, contents and linear
       recap; page 2 nonlinear model, derivative notation, rank condition and
       SSR expansion; page 3 residual-variance limit, FOCs, Taylor expansion
       and scaling; page 4 CLT, NLS limiting distribution and GNR covariance;
       page 5 parameter restrictions, power-model example and restricted GNR.
\end{transcriptnotes}

\end{document}
